\doublespacing % Do not change - required

\chapter{Related Work and Literature Review}
\label{ch2}

%%%%%%%%%%%%%%%%%%%%%%%%%%%%%%%%%%%%%%%
% IMPORTANT
\begin{spacing}{1} %THESE FOUR
\minitoc % LINES MUST APPEAR IN
\end{spacing} % EVERY
\thesisspacing % CHAPTER
% COPY THEM IN ANY NEW CHAPTER
%%%%%%%%%%%%%%%%%%%%%%%%%%%%%%%%%%%%%%%




\section{Dimensions of Post-Training Quantization}

LLM quantization reduces the numerical precision used to represent model tensors in order to lower storage, memory traffic and potentially inference cost. The practical difficulty increases as precision is reduced, particularly for Transformer models whose weight and activation distributions may contain sensitive channels or large-magnitude outliers. Quantization can therefore differ not only in bit width but also in scope, ranging from weight-only schemes to methods that additionally quantize runtime activations and the \ac{KV} cache. These distinctions provide the basis for the specific methods reviewed in the following sections.


\subsection{Post-Training Quantization and Calibration}

Post-training quantization converts a pretrained model into a
lower-precision representation without repeating the original
end-to-end training process. This distinguishes it from
quantization-aware training, in which simulated quantization effects are
introduced during optimisation so that the model parameters can adapt to
the resulting numerical error.
For large pretrained language models, avoiding full retraining makes
\ac{PTQ} attractive when the original training data or computational
budget are unavailable.

A typical \ac{PTQ} method uses representative calibration inputs to
estimate information needed to construct the low-precision model. The
required information is method dependent. It may include activation
ranges, clipping thresholds, channel statistics or layer-output
reconstruction targets. Calibration should therefore be regarded as part
of the quantizer construction rather than as final model evaluation.
The importance of calibration also means that results from different
published methods cannot be assumed to be directly comparable when they
use different calibration data, sample counts or numerical settings.


\subsection{Quantization Granularity}
For a conventional uniform scalar quantizer, a high-precision value
\(x\) is represented using an integer code and a scale factor. A generic
symmetric form may be written as
\begin{equation}
q(x)
=
\operatorname{clip}
\left(
\operatorname{round}\left(\frac{x}{s}\right),
q_{\min},q_{\max}
\right),
\qquad
\hat{x}=s\,q(x),
\label{eq:generic-scalar-quantizer}
\end{equation}
where \(s\) is the scale and \(\hat{x}\) is the reconstructed value.
Reducing the available bit width decreases the number of reconstruction
levels and therefore increases the difficulty of approximating the
original tensor accurately.

Quantization methods also differ in the granularity over which their
parameters are shared. Per-tensor quantization uses one scale for an
entire tensor, whereas per-channel quantization assigns separate scales
to individual channels. Group-wise quantization lies between these two
cases by assigning one scale to a smaller group of consecutive values.
Finer granularity can better adapt to local variations in numerical
range, but requires additional metadata and may complicate the runtime
implementation. ZeroQuant, for example, combines fine-grained weight
quantization with token-wise activation quantization to improve the
accuracy of low-precision Transformer inference.

\subsection{Scalar and Vector Representations}
Scalar quantization treats individual values independently, even when a
shared group scale is used. Vector quantization instead maps a block of
values jointly to a codeword from a multidimensional codebook. For
\(\mathbf{x}\in\mathbb{R}^{d}\) and codebook
\(\mathcal{C}\subset\mathbb{R}^{d}\), nearest-neighbour vector
quantization can be expressed as
\begin{equation}
Q_{\mathrm{V}}(\mathbf{x})
=
\underset{\mathbf{c}\in\mathcal{C}}
{\operatorname{arg\,min}}
\left\|
\mathbf{x}-\mathbf{c}
\right\|_{2}^{2}.
\label{eq:vector-quantization-ch2}
\end{equation}

The joint representation allows the codebook to exploit the geometry of
a block rather than allocating reconstruction levels independently to
each coordinate. Classical vector quantization therefore provides a
different rate--distortion trade-off from scalar rounding
\cite{1094577}. This distinction becomes important later in the chapter:
GPTQ, AWQ and SmoothQuant primarily operate with scalar or group-wise
representations, whereas QuIP\# and NestQuant introduce structured vector
codebooks.


\section{Weight-Only Post-Training Quantization}

Weight-only quantization reduces the precision of the fixed model
parameters while retaining runtime activations and the \ac{KV} cache at
higher precision. This setting is attractive because the weights
represent a large persistent component of the model and can be processed
offline. However, aggressive low-bit weight quantization can introduce
large layer-output errors, particularly when a small subset of channels
is more important to the computation than the remainder.

\subsection{\ac{GPTQ}}
GPTQ is a reconstruction-aware weight-only \ac{PTQ} method designed for
large Transformer models.

For a weight matrix \(\mathbf{W}\), calibration activations
\(\mathbf{X}\), and a set of admissible low-precision matrices
\(\mathcal{Q}\), the layer reconstruction problem may be written as
\begin{equation}
\widehat{\mathbf{W}}
=
\underset{\widetilde{\mathbf{W}}\in\mathcal{Q}}
{\operatorname{arg\,min}}
\left\|
\mathbf{W}\mathbf{X}
-
\widetilde{\mathbf{W}}\mathbf{X}
\right\|_{F}^{2}.
\label{eq:gptq-reconstruction-ch2}
\end{equation}

GPTQ uses approximate second-order information derived from the
calibration activations to compensate for the error introduced as
weights are quantized. Its block-wise implementation, shared ordering and
factorisation strategy make second-order reconstruction practical for
models containing billions of parameters. The method demonstrates that
preserving layer behaviour can be more informative than independently
minimising weight-space error.

The main limitation for the present dissertation is scope. GPTQ is
principally a weight-only method and does not, by itself, reduce
activation or KV-cache precision. It is therefore reviewed as an
important reconstruction-aware reference, but it is not one of the four
principal experimental routes evaluated in Chapter~\ref{ch4}.

\subsection{\ac{AWQ}}
\ac{AWQ} is a hardware-oriented weight-only post-training quantization
method developed for low-bit LLM inference. Its central observation is
that weight values are not equally important to model behaviour.
Experiments reported in the original work showed that protecting a
small fraction of salient weight channels can substantially reduce the
quality degradation caused by aggressive low-bit quantization.

AWQ determines salience using the distribution of the input
activations rather than the magnitude of the weights alone. Weight
channels receiving consistently large-magnitude activation features
have a greater influence on the corresponding linear-layer output and
are therefore treated as more sensitive to quantization. This is why
the method is described as \emph{activation-aware}, despite quantizing
only model weights.

One direct way to protect the salient channels would be to retain them
in FP16 while quantizing the remaining weights. Although this
mixed-precision strategy can preserve model quality, irregularly
interleaving floating-point and low-bit weights is difficult to execute
efficiently on conventional hardware. AWQ therefore protects salient
channels through an equivalent scaling transformation while retaining a
uniform low-bit weight representation.

For a linear layer with weight matrix \(\mathbf{W}\), input
activations \(\mathbf{X}\) and positive per-input-channel scale vector
\(\mathbf{s}\), the full-precision operation can be rewritten as

\begin{equation}
\mathbf{Y}
=
\mathbf{W}\mathbf{X}
=
\left(
\mathbf{W}\operatorname{diag}(\mathbf{s})
\right)
\left(
\operatorname{diag}(\mathbf{s})^{-1}\mathbf{X}
\right).
\label{eq:awq-equivalent-scaling}
\end{equation}

Equation~\eqref{eq:awq-equivalent-scaling} is mathematically equivalent
before quantization: scaling an input channel of the weight matrix can
be cancelled by inversely scaling the corresponding activation
channel. Once the transformed weights are quantized, however, scaling
up an important weight channel can reduce its relative rounding error.
The inverse activation scale may generally be fused into an adjacent
normalisation or preceding operation, avoiding an explicit additional
runtime multiplication.

AWQ selects the scaling vector by minimising the difference between the
full-precision and quantized layer outputs:

\begin{equation}
\mathbf{s}^{*}
=
\underset{\mathbf{s}}{\operatorname{arg\,min}}
\left\|
Q\left(
\mathbf{W}\operatorname{diag}(\mathbf{s})
\right)
\left(
\operatorname{diag}(\mathbf{s})^{-1}\mathbf{X}
\right)
-
\mathbf{W}\mathbf{X}
\right\|_{F}^{2},
\label{eq:awq-scaling-objective}
\end{equation}

where \(Q(\cdot)\) denotes the selected low-bit weight quantizer. The
search is constrained using the per-channel average activation
magnitude \(\mathbf{s}_{X}\), with candidate scales parameterised as

\begin{equation}
\mathbf{s}
=
\mathbf{s}_{X}^{\alpha},
\qquad
\alpha^{*}
=
\underset{\alpha\in[0,1]}{\operatorname{arg\,min}}
\mathcal{L}\left(\mathbf{s}_{X}^{\alpha}\right).
\label{eq:awq-scale-search}
\end{equation}

The scalar parameter \(\alpha\) controls the balance between protecting
salient and non-salient channels. An overly aggressive scale can reduce
the relative error of important channels while increasing the
quantization range and error experienced by other weights. AWQ
therefore performs a small grid search over \(\alpha\) and additionally
applies weight clipping to reduce quantization error.

Unlike GPTQ, AWQ does not sequentially reconstruct the weights using an
inverse Hessian and does not require backpropagation. The calibration
data are used primarily to estimate simple per-channel activation
statistics and to evaluate candidate scaling factors. The authors argue
that this reduced dependence on direct layer reconstruction limits
overfitting to the calibration samples and improves transfer across
tasks, domains and modalities. This represents an
important methodological contrast with second-order
reconstruction-based quantization.

The principal advantage of AWQ is that it provides an activation-aware
mechanism for protecting important features while retaining a regular
low-bit weight format. It avoids the irregularity of explicitly keeping
selected weights at higher precision and avoids the computational cost
of GPTQ-style second-order reconstruction. However, AWQ remains a
weight-only method: activations and the KV cache are retained at higher
precision in a conventional W4A16 configuration. Its benefits therefore
primarily concern static weight storage and weight-loading bandwidth,
rather than compression of every tensor involved in LLM inference.

\subsection{Critical comparison of weight-only approaches}

GPTQ and AWQ both use calibration information, but for different
purposes. GPTQ uses approximate second-order information to minimise
layer-output reconstruction error and compensate for sequential
rounding, whereas AWQ uses activation statistics to identify salient
channels and select scaling and clipping transformations that protect
them. GPTQ is therefore reconstruction-aware, while AWQ is primarily
activation-salience-aware.

Each approach also introduces distinct limitations. GPTQ performs
layer-wise reconstruction using calibration activations, which can make
its behaviour sensitive to the calibration distribution; AWQ reports
that reconstruction-based methods may overfit the calibration set. More recent work has
also identified error accumulation across independently calibrated
layers as a limitation of conventional GPTQ-style calibration and
proposed asymmetric calibration to compensate for this effect. AWQ avoids second-order reconstruction and
backpropagation, but remains fundamentally a weight-only method.
Consequently, activations and the KV cache remain at higher precision.
Moreover, practical W4A16 acceleration depends on efficient weight
packing, dequantization and specialised kernel support rather than on
the reduced weight bit width alone.

These limitations reinforce a broader point relevant to this
dissertation: a nominal W4 configuration does not uniquely determine
model quality or deployment behaviour. Quantizer design, calibration,
granularity, packing format, quantization scope and execution backend
must be considered together when comparing low-bit implementations.

\section{Weight--Activation Post-Training Quantization}

Weight--activation \ac{PTQ} extends low-precision representation from
the static model parameters to the dynamically generated activations
used during inference. Compared with weight-only quantization, this
broader scope can reduce memory traffic across a larger fraction of the
linear-layer data path and can enable integer or mixed-precision matrix
multiplication when suitable kernels are available. However, activation
quantization is more difficult because activation values depend on the
current input and may contain persistent large-magnitude outliers.
Methods such as ZeroQuant and SmoothQuant therefore introduce
fine-grained quantization, distribution transformation and specialised
runtime support to enable practical W8A8 inference. The following subsections review the activation outlier problem, SmoothQuant, the W4A8 operating regime and the
limitations of joint weight--activation quantization.

\subsection{Activation outliers}
Activation quantization presents a different challenge from weight
quantization. Model weights remain fixed after training and can
therefore be analysed, transformed and quantized offline. Activations,
by contrast, are generated dynamically during inference and depend on
the input sequence, token position and network layer. Their numerical
ranges are consequently not known exactly before deployment, requiring
activation scales either to be estimated from representative calibration
data or to be calculated dynamically at runtime.

A particularly important difficulty is the presence of
large-magnitude activation outliers. Consider a signed symmetric
\(b\)-bit quantizer applied to an activation group
\(\mathcal{G}\). Its scale may be defined as

\begin{equation}
s_{\mathcal{G}}
=
\frac{
\displaystyle\max_{x\in\mathcal{G}} |x|
}{
2^{b-1}-1
},
\label{eq:activation-outlier-scale}
\end{equation}

where \(2^{b-1}-1\) is the largest positive integer code available to
the quantizer. As indicated by
Equation~\eqref{eq:activation-outlier-scale}, the largest-magnitude
value within the selected group determines the spacing between adjacent
reconstruction levels. If only a small number of activation values are
substantially larger than the remainder of the distribution, the scale
must nevertheless be sufficiently large to represent them.

This enlarged scale reduces the effective resolution available to the
majority of non-outlier values. Values concentrated near zero occupy
only a small portion of the representable integer range, and multiple
distinct high-precision activations may consequently be mapped to the
same or neighbouring integer codes. Activation outliers can therefore
increase quantization error across an entire tensor or quantization
group even when the outlying values themselves constitute only a small
fraction of the activations.

\subsection{SmoothQuant}
SmoothQuant is a training-free W8A8 \ac{PTQ} method that addresses
activation outliers through an equivalent channel-wise transformation. Its central observation is that weights
are generally easier to quantize than activations. Quantization
difficulty can therefore be redistributed: activation channels are
scaled down while the corresponding weight channels are scaled up so
that the original floating-point linear transformation remains
unchanged before quantization.

The smoothing transformation is determined using calibration
statistics. A migration-strength parameter controls how much numerical
range is transferred from activations to weights. If too little
difficulty is transferred, activation outliers remain problematic; if
too much is transferred, the transformed weights themselves become
harder to quantize. SmoothQuant therefore treats weight and activation
quantization as a coupled problem rather than two independent
rounding operations.

An important systems advantage is that the smoothing transformation can
be folded into neighbouring model operations before inference. The
resulting model can then use a comparatively regular INT8 path. The
original work demonstrates that this approach can retain model quality
while reducing memory and improving inference efficiency across several
large language-model families. In this dissertation, SmoothQuant
therefore serves as the principal W8A8 reference. The exact smoothing
coefficient, activation granularity and execution backend used in the
project are intentionally deferred to Chapter~\ref{ch3}.

\subsection{W4A8-Related Approaches}
\ac{W4A8} denotes a mixed-precision configuration in which model
weights are represented using 4 bits and runtime activations using
8 bits. It is therefore a precision regime rather than a single
quantization algorithm. Conceptually, W4A8 combines the stronger
weight compression of W4A16 with the broader low-precision execution
scope of W8A8. However, reducing the weights from 8 bits to 4 bits
introduces substantially greater weight-quantization error, while
activation outliers must still be controlled during runtime.

FPTQ directly investigates post-training W4A8 quantization for large
language models. It combines fine-grained weight quantization with
layer-dependent activation strategies, including logarithmic
equalisation for layers that are particularly difficult to quantize. This work demonstrates that W4A8 generally requires
more specialised treatment than simply replacing the W8 weights of a
W8A8 model with a conventional 4-bit quantizer.

Other studies address the same precision regime through alternative
numerical representations. ZeroQuant-FP investigates FP4 weights and
FP8 activations, together with constrained weight scaling and
low-rank compensation, to reduce the degradation introduced by W4A8
quantization LQER instead combines activation-induced
scaling with a low-rank approximation of the quantization error,
allowing part of the error discarded by the 4-bit weight representation
to be reconstructed through an additional low-rank path.

These methods illustrate that the main challenge of W4A8 is not the
notation itself, but the coordination of weight quantization,
activation-range control and error reconstruction. Compared with
W8A8, W4A8 can reduce static weight storage further, but usually
requires finer weight granularity, improved scaling or an auxiliary
error-correction mechanism to maintain model quality. Compared with
W4A16, it additionally introduces input-dependent activation
quantization and its associated runtime overhead.

The systems benefit of W4A8 also depends on hardware and kernel support.
Four-bit weights may need to be unpacked, rescaled or converted before
they can be combined with 8-bit activations. Consequently, a W4A8
checkpoint can provide substantial storage compression without
automatically producing an equivalent latency improvement. Efficient
deployment requires an execution path capable of consuming the packed
weight representation while avoiding excessive quantization,
dequantization and precision-alignment overhead.

In this dissertation, W4A8 refers to the project-specific pipeline that
combines packed 4-bit weights with dynamically quantized 8-bit
activations. It is included as an intermediate operating point between
weight-only W4A16 quantization and the broader low-precision
configurations considered for NestQuant. The detailed calibration,
packing and execution procedures used by this implementation are
described in Chapter~\ref{ch3} rather than in the present literature review.

\subsection{Limitations of Joint Weight--Activation Quantization}
Joint weight--activation quantization introduces approximation error
into both the fixed model weights and the input-dependent activations,
making it inherently more challenging than weight-only quantization.
Activation distributions can contain persistent outliers and vary
between inputs, so calibration statistics cannot completely represent
every activation pattern encountered during deployment. SmoothQuant
addresses this problem by redistributing quantization difficulty from
activations to weights through channel-wise scaling, whereas
LLM.int8() retains selected outlier dimensions at higher precision.

However, these approaches introduce different compromises:
SmoothQuant remains dependent on representative calibration statistics
and an appropriate smoothing strength, while the mixed-precision
outlier path of LLM.int8() can introduce additional runtime overhead.
The SmoothQuant study itself shows that an unsuitable migration strength
can shift excessive quantization difficulty towards either activations
or weights, demonstrating that the two operands cannot be optimised
independently.

\section{Structured Vector and Lattice Quantization}
Scalar and group-wise quantizers allocate reconstruction values
independently to individual coordinates. Structured vector quantization
instead represents several values jointly, providing an opportunity to
use the geometry of the source distribution more efficiently. This has
motivated a recent line of LLM quantization work based on lattice and
other structured codebooks.

\subsection{Vector Quantization and Lattice Codebooks}
A lattice is a regular discrete set of points in a Euclidean vector
space. When used for quantization, the input is mapped to a nearby
lattice point and represented through the associated code. Structured
lattices can provide regular encoding and decoding procedures while
retaining the geometric benefits of vector quantization.

QuIP is an important example of this direction. It combines Hadamard
incoherence processing with lattice-based vector codebooks, including a
codebook derived from the eight-dimensional \(E_8\) lattice, to improve
very-low-bit weight quantization. This work demonstrates that structured vector
representations can outperform conventional scalar rounding in
weight-only compression. However, vector quantization also introduces
additional encoding and decoding complexity, which becomes especially
important when the quantizer is applied to tensors that must be processed
dynamically at runtime.

Recent information-theoretic work provides further motivation for
vector-codebook approaches. Under a Gaussian weight model, universal
codebooks can theoretically operate within a small rate gap of
codebooks adapted to the second-order statistics of the activations. The result is non-constructive and does
not itself provide an efficient deployment algorithm, but it reinforces
the idea that low-bit representation can be analysed as a
rate--distortion problem rather than only as scalar rounding.

\subsection{NestQuant}

NestQuant extends structured vector quantization to matrix products and
LLM inference using nested lattices. The practical
construction described in the original work uses the eight-dimensional
Gosset lattice \(E_8\), low-dimensional vector blocks, normalisation and
Hadamard rotation. Instead of using a single unrestricted codebook,
nested lattice structure provides a finite collection of efficiently
encodable reconstruction points.

The use of vector blocks provides a shaping advantage over independent
scalar quantization because bit patterns can be allocated according to a
multidimensional reconstruction region rather than a Cartesian grid.
The original NestQuant study argues that this becomes increasingly
important when the quantized vectors are approximately Gaussian after
normalisation and rotation.

A second distinguishing feature is quantization-aware LDLQ
(\mbox{QA-LDLQ}). Conventional reconstruction weight quantization
optimises the weights using statistics of the original activations.
When the activations themselves are also quantized, their quantization
noise changes the effective input seen by the quantized weight matrix.
QA-LDLQ modifies the weight-construction objective so that activation
quantization error is included during the reconstruction procedure. This makes NestQuant
particularly relevant to the broader W/A/KV quantization scope examined
in this dissertation.

NestQuant is also distinct from the other principal methods because its
scope is not limited to one tensor class. The framework can be applied to
weights, activations and the KV cache. This makes it suitable for the
controlled scope ablation performed later in the dissertation. The
specific \(q\), scaling-coefficient configuration, context length and
packed deployment path used in this project are implementation details
and are therefore reserved for Chapter~\ref{ch3}.

\subsection{Practical Boundaries of Structured Representations}
A vector quantization rate does not necessarily equal the complete
physical model footprint. Codeword indices, scale selections and
unquantized tensors may require additional storage. Likewise, a
numerically evaluated vector quantizer is not automatically a packed
runtime implementation. These distinctions are important when
interpreting effective bit rates or comparing vector quantization with
integer weight packing.

For this reason, the later NestQuant experiments separate two questions.
The A4 and KV4 configurations are used to study the sensitivity of model
quality to broader quantization scope, whereas physical checkpoint and
GPU-memory measurements are reported for the dedicated packed
W4A16KV16 deployment path. This avoids interpreting a numerical
quality result as evidence of an unmeasured runtime memory saving.

\section{KV-Cache Quantization}
Chapter~\ref{ch1} introduced the KV cache as a distinct component of
autoregressive inference memory. The literature on KV-cache quantization
therefore focuses on reducing the precision of the stored key and value
states while limiting disruption to the attention computation.

KIVI demonstrates that keys and values can exhibit different numerical
properties and therefore need not use the same quantization granularity.
It applies per-channel quantization to keys and per-token quantization to
values, enabling tuning-free low-bit cache storage. KVQuant similarly studies dedicated cache
quantization for long-context inference, including pre-RoPE key
quantization, non-uniform numerical representations and explicit
treatment of outliers. These approaches show that the KV cache is a
separate quantization target with different requirements from static
weight compression.

NestQuant provides another route by applying its vector quantizer to
keys and values in addition to weights and activations. This broader scope is
particularly useful for the present study because KV4 can be introduced
independently of A4. The resulting comparison allows the effect of
cache precision on model quality to be separated from the effect of
activation precision.

The literature nevertheless highlights an important boundary: a
lower-precision cache can only produce physical memory savings when the
cache is actually stored and processed in that compact representation.
The present dissertation therefore treats the NestQuant KV4 experiment
as a model-quality sensitivity study and does not claim packed KV4
runtime savings that were not directly measured.

\section{Research Gap}
The reviewed literature shows that low-bit LLM quantization is not a
single optimisation problem. GPTQ and AWQ primarily address low-bit
weights; SmoothQuant addresses joint weight--activation quantization;
W4A8-related methods introduce additional mechanisms for combining very
low-bit weights with runtime activation quantization; and NestQuant
extends structured vector quantization across weights, activations and
the KV cache. Dedicated methods such as KIVI and KVQuant further show
that cache quantization constitutes a separate design problem.

The methods also differ in calibration data, numerical formats,
quantization granularity, model families, evaluation datasets and
runtime backends. Published perplexity, memory or latency values
therefore cannot be treated as if they originated from one controlled
experiment. For the purposes of this dissertation, the relevant gap is
a practical comparative one: the selected low-bit operating points need
to be evaluated under a common model-quality protocol while keeping
quantization scope and physical representation distinct.

Table~\ref{tab:related-work-positioning} summarises how the principal
methods are positioned in the present study.

\begin{table}[htbp]
\centering
\caption{Positioning of the principal quantization approaches discussed
in this dissertation.}
\label{tab:related-work-positioning}

\begin{spacing}{1}
\small
\renewcommand{\arraystretch}{1.18}
\setlength{\tabcolsep}{4pt}

\begin{tabularx}{\textwidth}{@{}
>{\raggedright\arraybackslash}p{0.15\textwidth}
>{\raggedright\arraybackslash}X
>{\raggedright\arraybackslash}p{0.20\textwidth}
>{\raggedright\arraybackslash}X
@{}}
\toprule
\textbf{Approach} &
\textbf{Core principle} &
\textbf{Primary scope} &
\textbf{Role in this dissertation} \\
\midrule

GPTQ &
Second-order layer-output. &
Weights &
Methodological reference; not experimental baseline. \\

AWQ &
Activation-aware identification and scaling. &
Weights &
Principal weight-only route. \\

SmoothQuant &
Redistribution between activation and weight channels. &
Weights + activations &
Principal W8A8 route. \\

Project-specific W4A8 &
Packed low-bit weights and dynamic activation quantization. &
Weights + activations &
Mixed-precision operating point. \\

NestQuant &
Nested-lattice vector quantization with QA-LDLQ. &
Weights + activations + KV cache &
Quantization-scope ablation. \\

\bottomrule
\end{tabularx}
\end{spacing}
\end{table}

The dissertation therefore does not seek to establish that one published
algorithm is universally superior. Instead, it evaluates the retained
implementations according to four linked questions introduced in
Chapter~\ref{ch1}: model-quality retention, measured memory reduction,
sensitivity to activation and KV-cache quantization, and the physical
benefit of the packed NestQuant deployment path. Chapter~\ref{ch3}
specifies the exact model, calibration data, quantization pipelines and
measurement protocol used to answer these questions.

